Table~\ref{tab:results-main} compares Red-LGCP with the five baselines, with significance assessed by paired two-sided exact Wilcoxon signed-rank tests over the eight test weeks, Holm-corrected across baselines (App.~\ref{app:significance}). In both months, Red-LGCP achieves the best MAE, accuracy, F1, log-likelihood and CRPS, as well as the lowest Wasserstein distance among the predictive models, the lower value of LAP being an artifact of copying observed counts. These differences hold in all eight weeks and are significant (Holm-adjusted $p=0.039$, the smallest attainable value), with MAE reductions of 14--35\%; the only exception is CRPS against ConvLSTM, where Red-LGCP is better in seven weeks without reaching significance. The largest gains concern zeros: the covariate-based baselines predict a positive rate in every cell-day and thus score identically in accuracy and F1, whereas the classifier of Red-LGCP separates active from inactive cell-days, and its hurdle predictive distribution gives zeros their own probability, which also underlies its better log-likelihood and CRPS. ConvLSTM attains the lowest RMSE, which emphasizes rare large errors, but Red-LGCP is a close second in both months, and none of its RMSE differences from the predictive baselines is significant.
